\documentclass[10pt,a4paper]{amsart}
\usepackage{amsmath,amssymb,amsthm,graphicx}
\usepackage{verbatim}
\newenvironment{reference}{\comment}{\endcomment}
%\newenvironment{reference}{}{}
\newenvironment{slogan}{\comment}{\endcomment}
\newenvironment{history}{\comment}{\endcomment}

% For commutative diagrams we use Xy-pic
\usepackage[all]{xy}

% We use 2cell for 2-commutative diagrams.
\xyoption{2cell}
% Theorem environments.
%
\theoremstyle{plain}
\newtheorem{theorem}[subsection]{Theorem}
\newtheorem{proposition}[subsection]{Proposition}
\newtheorem{lemma}[subsection]{Lemma}

\theoremstyle{definition}
\newtheorem{definition}[subsection]{Definition}
\newtheorem{example}[subsection]{Example}
\newtheorem{exercise}[subsection]{Exercise}
\newtheorem{situation}[subsection]{Situation}

\theoremstyle{remark}
\newtheorem{remark}[subsection]{Remark}
\newtheorem{remarks}[subsection]{Remarks}

\numberwithin{equation}{subsection}

% Macros
%
\def\lim{\mathop{\mathrm{lim}}\nolimits}
\def\colim{\mathop{\mathrm{colim}}\nolimits}
\def\Spec{\mathop{\mathrm{Spec}}}
\def\Hom{\mathop{\mathrm{Hom}}\nolimits}
\def\Ext{\mathop{\mathrm{Ext}}\nolimits}
\def\SheafHom{\mathop{\mathcal{H}\!\mathit{om}}\nolimits}
\def\SheafExt{\mathop{\mathcal{E}\!\mathit{xt}}\nolimits}
\def\Sch{\mathit{Sch}}
\def\Mor{\mathop{\mathrm{Mor}}\nolimits}
\def\Ob{\mathop{\mathrm{Ob}}\nolimits}
\def\Sh{\mathop{\mathit{Sh}}\nolimits}
\def\NL{\mathop{N\!L}\nolimits}
\def\CH{\mathop{\mathrm{CH}}\nolimits}
\def\proetale{{pro\text{-}\acute{e}tale}}
\def\etale{{\acute{e}tale}}
\def\QCoh{\mathit{QCoh}}
\def\Ker{\mathop{\mathrm{Ker}}}
\def\Im{\mathop{\mathrm{Im}}}
\def\Coker{\mathop{\mathrm{Coker}}}
\def\Coim{\mathop{\mathrm{Coim}}}

% Boxtimes
%
\DeclareMathSymbol{\boxtimes}{\mathbin}{AMSa}{"02}

%
% Macros for moduli stacks/spaces
%
\def\QCohstack{\mathcal{QC}\!\mathit{oh}}
\def\Cohstack{\mathcal{C}\!\mathit{oh}}
\def\Spacesstack{\mathcal{S}\!\mathit{paces}}
\def\Quotfunctor{\mathrm{Quot}}
\def\Hilbfunctor{\mathrm{Hilb}}
\def\Curvesstack{\mathcal{C}\!\mathit{urves}}
\def\Polarizedstack{\mathcal{P}\!\mathit{olarized}}
\def\Complexesstack{\mathcal{C}\!\mathit{omplexes}}
% \Pic is the operator that assigns to X its picard group, usage \Pic(X)
% \Picardstack_{X/B} denotes the Picard stack of X over B
% \Picardfunctor_{X/B} denotes the Picard functor of X over B
\def\Pic{\mathop{\mathrm{Pic}}\nolimits}
\def\Picardstack{\mathcal{P}\!\mathit{ic}}
\def\Picardfunctor{\mathrm{Pic}}
\def\Deformationcategory{\mathcal{D}\!\mathit{ef}}
\usepackage{cleveref}
\crefname{originallemma}{original lemma}{original lemmas}
\Crefname{originallemma}{Original lemma}{Original lemmas}
\makeatletter
\newlabel{algebra-lemma-completion-Noetherian}{{05GH}{source}{}{sourcepointer-1}{}}
\newlabel{algebra-lemma-completion-Noetherian@cref}{{[originallemma][0][]05GH}{[1][0][]source}{}{}{}}
\makeatother
\crefname{originalproposition}{original proposition}{original propositions}
\Crefname{originalproposition}{Original proposition}{Original propositions}
\makeatletter
\newlabel{algebra-proposition-dimension}{{00KQ}{source}{}{sourcepointer-2}{}}
\newlabel{algebra-proposition-dimension@cref}{{[originalproposition][0][]00KQ}{[1][0][]source}{}{}{}}
\makeatother
\title{The limit of the sections of the thickenings of an exceptional curve of the first kind is a regular local ring of dimension two}
\author{The Stacks Project Authors}
\date{stacks-project:resolve.tex:local-snapshot-20261012; GNU FDL 1.2 or later}
\begin{document}
\maketitle
\noindent\textit{Editorial note.} All delivered mathematical text is the author's own, copied from the pinned Stacks Project chapter file resolve.tex; no statement or proof step was rewritten, reordered or removed.\par\medskip
\noindent\textit{Editorial note.} The counted proof is the author's complete argument: the use of the vanishing of the first cohomology of the twists to obtain surjective transition maps and the description of the maximal ideal of the limit, the choice of two generators whose monomials give bases of the successive quotients, and the expansion argument identifying each kernel with a power of the ideal generated by those two elements.\par\medskip
\noindent\textit{Editorial note.} The label lemma-exceptional-first-kind-local names the counted statement, and the proof environment immediately following it in the source is delivered verbatim as its proof.\par\medskip
\noindent\textit{Editorial note.} The counted statement is carried by the author's own lemma environment in the chapter source.\par\medskip
\noindent\textit{Editorial note.} Pre-open gate: python operations/agent-authoring/tools/precard\_gate.py --source resolve.tex --label lemma-exceptional-first-kind-local --start 4603 --end 4675 --local-id 28078 returned RESULT=OK (L28078 free, occupancy FREE, public interval clear). Pre-assembly tag self-check: canonical tag resolved and not present in reference/public-items.jsonl.\par\medskip
\noindent\textit{Editorial note.} The statement and proof are copied from the pinned chapter file; the cited inputs are kept as external pointers with their canonical Stacks tags.\par\medskip
\noindent\textit{Allowed standard prerequisites (editorial).}\begin{itemize}
\item Description of an exceptional curve of the first kind. Same chapter, the definition of an exceptional curve of the first kind as a curve isomorphic to a projective line whose normal sheaf pulls back to the twist by minus one. Supplies the identification of the curve with a projective line and the computation of the cohomology of the twists by means of the normal sheaf.
\item The ideal sheaf of a thickening and the associated short exact sequences. Standard description of the ideal sheaf of the multiple nE of an effective Cartier divisor as the invertible sheaf of the negative multiple, giving the displayed short exact sequences. Supplies the exact sequences comparing the sections over successive thickenings with the cohomology of the curve itself.
\item Noetherianness and dimension of an adically complete local ring. Algebra, the lemma that the completion of a Noetherian local ring is Noetherian and the proposition expressing the dimension through the growth of the powers of the maximal ideal. Supply the Noetherianness and the dimension count for the limit ring.
\end{itemize}\medskip
\ifdefined\hypertarget\hypertarget{sourcepointer-1}{}\fi
\noindent\textit{Source reference.} Original lemma 05GH: sources/stacks-chow/tags.txt, source lines 7061--7061. Prerequisite pointer invoked at proof line(s) 4669; it names a standard fact of a sibling chapter that the counted argument uses, and every step of the counted proof built on it is retained verbatim in the delivered body.\par\medskip
\ifdefined\hypertarget\hypertarget{sourcepointer-2}{}\fi
\noindent\textit{Source reference.} Original proposition 00KQ: sources/stacks-chow/tags.txt, source lines 760--760. Prerequisite pointer invoked at proof line(s) 4672; it names a standard fact of a sibling chapter that the counted argument uses, and every step of the counted proof built on it is retained verbatim in the delivered body.\par\medskip
\begin{lemma}
\label{lemma-exceptional-first-kind-local}
Let $X$ be a Noetherian scheme. Let $E \subset X$ be an
exceptional curve of the first kind. Let $E_n = nE$ and
denote $\mathcal{O}_n$ its structure sheaf. Then
$$
A = \lim H^0(E_n, \mathcal{O}_n)
$$
is a complete local Noetherian regular local ring of dimension $2$
and $\Ker(A \to H^0(E_n, \mathcal{O}_n))$ is the $n$th power of
its maximal ideal.
\end{lemma}

\begin{proof}
Recall that there exists an isomorphism $\mathbf{P}^1_k \to E$
such that the normal sheaf of $E$ in $X$ pulls back to $\mathcal{O}(-1)$.
Then $H^0(E, \mathcal{O}_E) = k$.
We will denote $\mathcal{O}_n(iE)$ the restriction of the invertible
sheaf $\mathcal{O}_X(iE)$ to $E_n$ for all $n \geq 1$ and $i \in \mathbf{Z}$.
Recall that $\mathcal{O}_X(-nE)$ is the ideal sheaf of $E_n$. Hence
for $d \geq 0$ we obtain a short exact sequence
$$
0 \to \mathcal{O}_E(-(d + n)E) \to
\mathcal{O}_{n + 1}(-dE) \to
\mathcal{O}_n(-dE) \to 0
$$
Since $\mathcal{O}_E(-(d + n)E) = \mathcal{O}_{\mathbf{P}^1_k}(d + n)$
the first cohomology group vanishes for all $d \geq 0$ and $n \geq 1$.
We conclude that the transition maps of the system
$H^0(E_n, \mathcal{O}_n(-dE))$ are surjective. For $d = 0$
we get an inverse system of surjections of rings such that the
kernel of each transition map is a nilpotent ideal.
Hence $A = \lim H^0(E_n, \mathcal{O}_n)$ is a local ring
with residue field $k$ and maximal ideal
$$
\lim \Ker(H^0(E_n, \mathcal{O}_n) \to H^0(E, \mathcal{O}_E)) =
\lim H^0(E_n, \mathcal{O}_n(-E))
$$
Pick $x, y$ in this kernel mapping to a $k$-basis of
$H^0(E, \mathcal{O}_E(-E)) = H^0(\mathbf{P}^1_k, \mathcal{O}(1))$.
Then $x^d, x^{d - 1}y, \ldots, y^d$ are
elements of $\lim H^0(E_n, \mathcal{O}_n(-dE))$ which map to a basis
of $H^0(E, \mathcal{O}_E(-dE)) = H^0(\mathbf{P}^1_k, \mathcal{O}(d))$.
In this way we see that $A$ is separated and complete with respect
to the linear topology defined by the kernels
$$
I_n = \Ker(A \longrightarrow H^0(E_n, \mathcal{O}_n))
$$
We have $x, y \in I_1$, $I_d I_{d'} \subset I_{d + d'}$
and $I_d/I_{d + 1}$ is a free $k$-module on $x^d, x^{d - 1}y, \ldots, y^d$.
We will show that $I_d = (x, y)^d$. Namely, if $z_e \in I_e$ with
$e \geq d$, then we can write
$$
z_e = a_{e, 0} x^d + a_{e, 1} x^{d - 1}y + \ldots + a_{e, d}y^d + z_{e + 1}
$$
where $a_{e, j} \in (x, y)^{e - d}$ and $z_{e + 1} \in I_{e + 1}$
by our description of $I_d/I_{d + 1}$. Thus starting with some
$z = z_d \in I_d$ we can do this inductively
$$
z = \sum\nolimits_{e \geq d} \sum\nolimits_j a_{e, j} x^{d - j} y^j
$$
with some $a_{e, j} \in (x, y)^{e - d}$. Then $a_j = \sum_{e \geq d} a_{e, j}$
exists (by completeness and the fact that $a_{e, j} \in I_{e - d}$)
and we have $z = \sum a_{e, j} x^{d - j} y^j$.
Hence $I_d = (x, y)^d$.
Thus $A$ is $(x, y)$-adically complete. Then $A$ is
Noetherian by Algebra, Lemma \ref{algebra-lemma-completion-Noetherian}.
It is clear that the dimension is $2$ by the description
of $(x, y)^d/(x, y)^{d + 1}$ and
Algebra, Proposition \ref{algebra-proposition-dimension}.
Since the maximal ideal
is generated by two elements it is regular.
\end{proof}

\end{document}
